% Generated draft: requires mathematical and source review.
\documentclass[11pt]{article}
\usepackage[margin=25mm]{geometry}
\usepackage{fontspec}
\setmainfont{Latin Modern Roman}
\usepackage{amsmath,amssymb,mathtools,unicode-math}
\setmathfont{Latin Modern Math}
\usepackage{enumitem,needspace,xurl,hyperref}
\setlength{\parindent}{0pt}
\setlength{\parskip}{5pt}
\setlength{\emergencystretch}{3em}
\newcommand{\N}{\mathbb{N}}
\newcommand{\Z}{\mathbb{Z}}
\newcommand{\Q}{\mathbb{Q}}
\newcommand{\R}{\mathbb{R}}
\newcommand{\C}{\mathbb{C}}
\newcommand{\F}{\mathbb{F}}
\newcommand{\A}{\mathbb{A}}
\newcommand{\PP}{\mathbb{P}}
\newcommand{\E}{\mathbb{E}}
\newcommand{\eps}{\varepsilon}
\newcommand{\dd}{\,\mathrm{d}}
\begin{document}
\section*{UnsolvedMath ID: 800284}
\textbf{Source problem code:} catalogue entry
\subsection{Definitions and mathematical statement}
Let \(n\geq 1\) and let \(\Omega\subseteq\C^n\) be a bounded domain, meaning a connected open set. Write \(d_\Omega\) for its Kobayashi distance and assume that \((\Omega,d_\Omega)\) is complete. For \(z\in\Omega\) and \(r>0\), define the Kobayashi ball
\[
B_\Omega(z,r)=\{\zeta\in\Omega:d_\Omega(z,\zeta)<r\}.
\]
Let \(\mathcal O(\Omega)\) denote the holomorphic functions on \(\Omega\). For the Euclidean-volume interpretation of the source, let \(dV\) be Lebesgue measure on \(\C^n\cong\R^{2n}\) and let \(|B_\Omega(z,r)|=\int_{B_\Omega(z,r)}dV\).

The source asks whether, for every \(p>0\), there is a constant \(C=C(p)>0\) such that, for every \(f\in\mathcal O(\Omega)\), every \(z\in\Omega\), and every sufficiently small \(r>0\),
\[
|f(z)|^p\leq \frac{C}{|B_\Omega(z,r)|}\int_{B_\Omega(z,r)}|f(\zeta)|^p\,dV(\zeta).
\]
The constant is intended to be uniform in \(f\), \(z\), and the permitted radii. However, the source does not specify whether the small-radius threshold is uniform in \(z\) and \(f\), or what other parameters it may depend on. Although the notation \(C(p)\) lists only \(p\), the domain is fixed at the outset, so it is also unclear whether dependence on \(\Omega\) or \(n\) is permitted.

The measure is explicitly left open in the source. Besides Euclidean volume, it permits consideration of canonical measures associated with the Kobayashi construction, mentioning measures due to Eisenman or Bun Wong. For a chosen such measure \(\mu\), the corresponding normalized inequality would be
\[
|f(z)|^p\leq \frac{C}{\mu(B_\Omega(z,r))}\int_{B_\Omega(z,r)}|f(\zeta)|^p\,d\mu(\zeta),
\]
where the normalization requires \(0<\mu(B_\Omega(z,r))<\infty\). The source does not specify a canonical measure or whether the question asks for at least one admissible measure or for every suggested measure. Thus these measure choices and uniformity requirements remain unresolved parts of the formulation.

As background, the source reports that the result holds when \(\Omega\) is of finite type in \(\C^2\), or is strongly pseudoconvex in any dimension, citing comparability of metric balls with polydiscs and classical arguments of Stein. This is a supplied assertion, not an independently verified theorem here. The passage does not specify boundary regularity, the meaning of finite type, or the precise hypotheses and estimates used in that assertion.
\subsection{Short English statement}
For a bounded domain with complete Kobayashi distance, ask whether the value of \(|f|^p\) at the center of every sufficiently small Kobayashi ball is bounded by a uniform constant times its average over that ball. The source leaves the measure, the permitted dependence of the constant, and the uniformity of the radius threshold unresolved. It reports positive special cases for finite-type domains in \(\C^2\) and strongly pseudoconvex domains. The supplied registry labels the problem open; this draft does not independently verify that status.
\subsection{Sources}
\begin{itemize}
\item Local UnsolvedMath registry, supplied complete record: permanent ID \(800284\), exact source problem code \texttt{catalogue entry}, title ``Several Complex Variables Problem List --- Problem 84.'' The record lists Buzzard, Fornaess, and Krantz as proposers, gives the proposed year as 1995, and labels the problem \texttt{open}. These are supplied registry attributes.
\item Buzzard, Fornaess, and Krantz, \emph{The Several Complex Variables Problem List} (1995), Problem 84. Supplied source URL: \url{https://arxiv.org/abs/math/9509202}. The record reports access on 2026-10-16 and extraction by \texttt{arxiv\_latex}; these are supplied provenance details, not a claim of independent access or verification.
\item The supplied problem passage is signed S. Krantz, with contact \texttt{sk@math.wustl.edu}. Its special-case discussion refers to classical arguments of Stein and to the contributor's ``SCV book,'' without supplying bibliographic details. It also mentions Eisenman and Bun Wong in connection with possible canonical measures, without identifying particular publications or constructions.
\end{itemize}
\end{document}
